
\FigureLayoutDeclare{img/2018/jiangsu/q18_fig1.png}{6.0cm}{193bp 184bp 177bp 0bp}

\chapter{2018年江苏卷}

\section{填空题}

\begin{problem}
已知集合 \(A=\{0,1,2,8\}\)，\(B=\{-1,1,6,8\}\)，那么 \(A\cap B=\)\fillinblank{}.
\end{problem}

\begin{answer}
\(\{1,8\}\)
\end{answer}

\begin{solution}
交集由同时属于 \(A\) 与 \(B\) 的元素组成. \(A=\{0,1,2,8\}\) 与 \(B=\{-1,1,6,8\}\) 的公共元素为 \(1\)，\(8\)，所以 \(A\cap B=\{1,8\}\).
\end{solution}

\begin{problem}
若复数 \(z\) 满足 \(\i z=1+2\i\)，其中 \(\i\) 是虚数单位，则 \(z\) 的实部为\fillinblank{}.
\end{problem}

\begin{answer}
\(2\)
\end{answer}

\begin{solution}
由 \(\frac{1}{\i}=-\i\) 得 \(z=\frac{1+2\i}{\i}=(1+2\i)(-\i)=2-\i\)，其实部为 \(2\).
\end{solution}

\begin{problem}
已知 5 位裁判给某运动员打出的分数的茎叶图如图所示，那么这 5 位裁判打出的分数的平均数为\fillinblank{}.

\begin{center}
\begin{tabular}{r|lll}
8 & 9 & 9 & \\
9 & 0 & 1 & 1
\end{tabular}
\end{center}
\end{problem}

\begin{answer}
\(90\)
\end{answer}

\begin{solution}
茎叶图给出的五个分数为 \(89\)，\(89\)，\(90\)，\(91\)，\(91\)，其平均数为 \(\frac{89+89+90+91+91}{5}=\frac{450}{5}=90\).
\end{solution}

\begin{problem}
一个算法的伪代码如图所示，执行此算法，最后输出的 \(S\) 的值为\fillinblank{}.

\begin{center}
\bitmapfigure[width=6.0cm]{img/2018/jiangsu/q4_fig1.png}
\end{center}
\end{problem}

\begin{answer}
\(8\)
\end{answer}

\begin{solution}
初始 \(I=1\)，\(S=1\). 当 \(I<6\) 时执行 \(I\leftarrow I+2\)，\(S\leftarrow2S\)，依次得 \((I,S)=(3,2)\)，\((5,4)\)，\((7,8)\)；此时 \(I=7\) 不再满足循环条件，输出 \(S=8\).
\end{solution}

\begin{problem}
函数 \(f(x)=\sqrt{\log_2 x-1}\) 的定义域为\fillinblank{}.
\end{problem}

\begin{answer}
\([2,+\infty)\)
\end{answer}

\begin{solution}
被开方数非负要求 \(\log_2x-1\geqslant0\)，即 \(\log_2x\geqslant1\)，又真数 \(x>0\)，故 \(x\geqslant2\)，定义域为 \([2,+\infty)\).
\end{solution}

\begin{problem}
某兴趣小组有 2 名男生和 3 名女生，现从中任选 2 名学生去参加活动，则恰好选中 2 名女生的概率为\fillinblank{}.
\end{problem}

\begin{answer}
\(\frac{3}{10}\)
\end{answer}

\begin{solution}
从 \(5\) 人中任选 \(2\) 人有 \(\mathrm C_5^2=10\) 种等可能结果，其中 \(2\) 人均为女生的有 \(\mathrm C_3^2=3\) 种，所求概率为 \(\frac{3}{10}\).
\end{solution}

\begin{problem}
已知函数 \(y=\sin(2x+\varphi)\)（\(-\dfrac{\pi}{2}<\varphi<\dfrac{\pi}{2}\)） 的图象关于直线 \(x=\dfrac{\pi}{3}\) 对称，则 \(\varphi\) 的值是\fillinblank{}.
\end{problem}

\begin{answer}
\(-\frac{\pi}{6}\)
\end{answer}

\begin{solution}
图象关于直线 \(x=\frac{\pi}{3}\) 对称，故 \(2\cdot\frac{\pi}{3}+\varphi=\frac{\pi}{2}+k\pi\)（\(k\in\Z\)），即 \(\varphi=-\frac{\pi}{6}+k\pi\). 由 \(-\frac{\pi}{2}<\varphi<\frac{\pi}{2}\) 得 \(k=0\)，\(\varphi=-\frac{\pi}{6}\).
\end{solution}

\begin{problem}
在平面直角坐标系 \(xOy\) 中，若双曲线 \(\dfrac{x^2}{a^2}-\dfrac{y^2}{b^2}=1\)（\(a>0\)，\(b>0\)）的右焦点 \(F(c,0)\) 到一条渐近线的距离为 \(\dfrac{\sqrt3}{2}c\)，则其离心率的值是\fillinblank{}.
\end{problem}

\begin{answer}
\(2\)
\end{answer}

\begin{solution}
右焦点 \(F(c,0)\) 到渐近线 \(bx-ay=0\) 的距离为 \(\frac{bc}{\sqrt{a^2+b^2}}=\frac{bc}{c}=b\). 由题意 \(b=\frac{\sqrt3}{2}c\)，故 \(a^2=c^2-b^2=\frac{c^2}{4}\)，离心率 \(e=\frac{c}{a}=2\).
\end{solution}

\begin{problem}
函数 \(f(x)\) 满足 \(f(x+4)=f(x)\)（\(x\in\R\)），且在区间 \((-2,2]\) 上，
\(f(x)=\begin{cases}
\cos\dfrac{\pi x}{2}, & 0<x\le2,\\
\left|x+\dfrac12\right|, & -2<x\le0,
\end{cases}\)
则 \(f(f(15))\) 的值为\fillinblank{}.
\end{problem}

\begin{answer}
\(\frac{\sqrt{2}}{2}\)
\end{answer}

\begin{solution}
由周期 \(4\) 得 \(f(15)=f(15-4\cdot4)=f(-1)\). \(-1\in(-2,0]\)，故 \(f(-1)=\left|-1+\frac12\right|=\frac12\). 又 \(\frac12\in(0,2]\)，故 \(f(f(15))=f\left(\frac12\right)=\cos\frac{\pi}{4}=\frac{\sqrt2}{2}\).
\end{solution}

\begin{problem}
如图所示，正方体的棱长为 \(2\)，以其所有面的中心为顶点的多面体的体积为\fillinblank{}.

\begin{center}
\bitmapfigure[width=6.0cm]{img/2018/jiangsu/q10_fig1.png}
\end{center}
\end{problem}

\begin{answer}
\(\frac{4}{3}\)
\end{answer}

\begin{solution}
六个面心连成的多面体关于正方体中心对称；过中心的水平截面为正方形，其边长为相邻面心距离 \(\sqrt{1^2+1^2}=\sqrt2\)，面积为 \(2\)，上下两顶点到该截面的距离均为 \(1\). 故体积为 \(2\cdot\frac13\cdot2\cdot1=\frac43\).
\end{solution}

\begin{problem}
若函数 \(f(x)=2x^3-ax^2+1\)（\(a\in\R\)）在 \((0,+\infty)\) 内有且只有一个零点，则 \(f(x)\) 在 \([-1,1]\) 上的最大值与最小值的和为\fillinblank{}.
\end{problem}

\begin{answer}
\(-3\)
\end{answer}

\begin{solution}
\(f'(x)=6x^2-2ax=2x(3x-a)\). 若 \(a\leqslant0\)，则 \(f'(x)>0\) 在 \((0,+\infty)\) 上恒成立，\(f\) 严格递增，又 \(f(0)=1>0\)，故 \(f\) 在 \((0,+\infty)\) 内无零点，不合要求. 于是 \(a>0\)，\(f\) 在 \(\left(0,\frac{a}{3}\right)\) 上递减，在 \(\left(\frac{a}{3},+\infty\right)\) 上递增；有且仅有一个零点当且仅当 \(f\left(\frac{a}{3}\right)=2\left(\frac{a}{3}\right)^3-a\left(\frac{a}{3}\right)^2+1=-\frac{a^3}{27}+1=0\)，即 \(a=3\). 此时 \(f(x)=2x^3-3x^2+1\)，\(f'(x)=6x(x-1)\)，故 \(f\) 在 \([-1,0]\) 上递增，在 \([0,1]\) 上递减，最大值为 \(f(0)=1\)，最小值为 \(\min\{f(-1),f(1)\}=\min\{-4,0\}=-4\)，和为 \(-3\).
\end{solution}

\begin{problem}
在平面直角坐标系 \(xOy\) 中，\(A\) 为直线 \(l:y=2x\) 上在第一象限内的点，\(B(5,0)\)，以 \(AB\) 为直径的圆 \(C\) 与直线 \(l\) 交于另一点 \(D\). 若 \(\overrightarrow{AB}\cdot\overrightarrow{CD}=0\)，则点 \(A\) 的横坐标为\fillinblank{}.

\bitmapfigure[width=0.42\linewidth]{img/2018/jiangsu/q12_fig1.png}
\end{problem}

\begin{answer}
\(3\)
\end{answer}

\begin{solution}
设 \(A=(t,2t)\)（\(t>0\)）. 以 \(AB\) 为直径的圆方程为 \((x-t)(x-5)+(y-2t)y=0\)，与 \(y=2x\) 联立得 \((x-t)(5x-5)=0\)，除 \(A\) 外另一交点为 \(D=(1,2)\). 圆心为 \(C=\left(\frac{t+5}{2},t\right)\)，故 \(\overrightarrow{AB}=(5-t,-2t)\)，\(\overrightarrow{CD}=\left(1-\frac{t+5}{2},2-t\right)\). 由 \(\overrightarrow{AB}\cdot\overrightarrow{CD}=0\) 得 \(\frac{(5-t)(-t-3)}{2}-2t(2-t)=0\)，即 \((t-3)(t+1)=0\)，结合 \(t>0\) 得 \(t=3\).
\end{solution}

\begin{problem}
在 \(\triangle ABC\) 中，角 \(A\)，\(B\)，\(C\) 所对的边分别为 \(a\)，\(b\)，\(c\)，\(\angle ABC=120^\circ\)，\(\angle ABC\) 的平分线交 \(AC\) 于点 \(D\)，且 \(BD=1\)，则 \(4a+c\) 的最小值为\fillinblank{}.
\end{problem}

\begin{answer}
\(9\)
\end{answer}

\begin{solution}
记 \(BC=a\)，\(AB=c\). 由 \(S_{\triangle ABC}=S_{\triangle ABD}+S_{\triangle CBD}\) 得 \(\frac12ac\sin120^{\circ}=\frac12c\sin60^{\circ}+\frac12a\sin60^{\circ}\)，即 \(ac=a+c\)，亦即 \(\frac1a+\frac1c=1\). 于是 \(4a+c=(4a+c)\left(\frac1a+\frac1c\right)=5+\frac{4a}{c}+\frac{c}{a}\geqslant5+2\sqrt4=9\)，等号在 \(\frac{4a}{c}=\frac{c}{a}\) 即 \(c=2a\)（此时 \(a=\frac32\)，\(c=3\)）时成立，故最小值为 \(9\).
\end{solution}

\begin{problem}
已知集合 \(A=\{x\mid x=2n-1,\ n\in\N^*\}\)，\(B=\{x\mid x=2^n,\ n\in\N^*\}\). 将 \(A\cup B\) 的所有元素从小到大依次排列构成一个数列 \(\{a_n\}\). 记 \(S_n\) 为数列 \(\{a_n\}\) 的前 \(n\) 项和，则使得 \(S_n>12a_{n+1}\) 成立的 \(n\) 的最小值为\fillinblank{}.
\end{problem}

\begin{answer}
\(27\)
\end{answer}

\begin{solution}
若 \(a_{n+1}=2k+1\) 为奇数，则前 \(n\) 项为奇数 \(1,3,\ldots,2k-1\)（\(k\) 个）及小于 \(2k+1\) 的 \(2\) 的幂 \(2,\ldots,2^m\)（\(2^m<2k+1\)，\(m\) 个），故 \(S_n=k^2+(2^{m+1}-2)\)，不等式化为 \(k^2-24k-14+2^{m+1}>0\). 当 \(m\leqslant4\) 时 \(k\leqslant15\)，左端不大于 \(k^2-24k+18<0\)；当 \(m=5\) 时化为 \(k^2-24k+50>0\)，得 \(k\geqslant22\)，此时 \(n=k+m\geqslant27\)；当 \(m\geqslant6\) 时 \(k\geqslant32\)，不等式恒成立，此时 \(n>27\). 若 \(a_{n+1}=2^m\)，则 \(n=2^{m-1}+m-1\)，\(S_n=2^{2m-2}+2^{m+1}-2\)，不等式化为 \(2^{2m-2}-10\cdot2^m-2>0\)，\(m\leqslant5\) 时不成立，\(m\geqslant6\) 时 \(n\geqslant37\). 直接验算 \(a_{28}=45\)，\(S_{27}=22^2+(2+4+8+16+32)=546>540=12\cdot45\). 故成立的 \(n\) 的最小值为 \(27\).
\end{solution}

\section{解答题}

\begin{problem}
在平行六面体 \(ABCD-A_1B_1C_1D_1\) 中，\(AA_1=AB\)，\(AB_1\perp B_1C_1\). 求证：

\begin{enumerate}
\item \(AB\parallel\) 平面 \(A_1B_1C\)；
\item 平面 \(ABB_1A_1\perp\) 平面 \(A_1BC\).
\end{enumerate}

\begin{center}
\bitmapfigure[width=0.42\linewidth]{img/2018/jiangsu/q15_fig2.png}
\end{center}
\end{problem}

\begin{solution}
\begin{enumerate}
\item 平行六面体的侧面 \(ABB_1A_1\) 为平行四边形，故 \(AB\parallel A_1B_1\). \(A_1B_1\subset\) 平面 \(A_1B_1C\) 而 \(AB\) 在该平面外，所以 \(AB\parallel\) 平面 \(A_1B_1C\).
\item 由 \(AA_1=AB\) 知 \(ABB_1A_1\) 为菱形，故 \(AB_1\perp A_1B\). 又 \(AB_1\perp B_1C_1\) 且 \(B_1C_1\parallel BC\)，得 \(AB_1\perp BC\). \(A_1B\) 与 \(BC\) 相交于 \(B\)，于是 \(AB_1\perp\) 平面 \(A_1BC\)；而 \(AB_1\subset\) 平面 \(ABB_1A_1\)，故平面 \(ABB_1A_1\perp\) 平面 \(A_1BC\).
\end{enumerate}
\end{solution}

\begin{problem}
已知 \(\alpha\)，\(\beta\) 为锐角，\(\tan\alpha=\dfrac{4}{3}\)，\(\cos(\alpha+\beta)=-\dfrac{\sqrt5}{5}\).

\begin{enumerate}
\item 求 \(\cos 2\alpha\) 的值；
\item 求 \(\tan(\alpha-\beta)\) 的值.
\end{enumerate}
\end{problem}

\begin{solution}
\begin{enumerate}
\item 由万能公式 \(\cos2\alpha=\frac{1-\tan^2\alpha}{1+\tan^2\alpha}=\frac{1-\frac{16}{9}}{1+\frac{16}{9}}=-\frac7{25}\).
\item \(\alpha\)，\(\beta\) 为锐角，故 \(\alpha+\beta\in(0,\pi)\)；由 \(\cos(\alpha+\beta)=-\frac{\sqrt5}{5}\) 得 \(\sin(\alpha+\beta)=\frac{2\sqrt5}{5}\)，\(\tan(\alpha+\beta)=-2\). 又 \(\tan2\alpha=\frac{2\tan\alpha}{1-\tan^2\alpha}=-\frac{24}{7}\)，而 \(\alpha-\beta=2\alpha-(\alpha+\beta)\)，故 \(\tan(\alpha-\beta)=\frac{\tan2\alpha-\tan(\alpha+\beta)}{1+\tan2\alpha\tan(\alpha+\beta)}=\frac{-\frac{24}{7}+2}{1+\frac{48}{7}}=-\frac2{11}\).
\end{enumerate}
\end{solution}

\begin{problem}
某农场有一块农田，如图所示，它的边界由圆 \(O\) 的一段圆弧 \(MPN\)（\(P\) 为此圆弧的中点）和线段 \(MN\) 构成. 已知圆 \(O\) 的半径为 40 米，点 \(P\) 到 \(MN\) 的距离为 50 米. 现规划在此农田上修建两个温室大棚，大棚 \(\mathrm{I}\) 内的地块形状为矩形 \(ABCD\)，大棚 \(\mathrm{II}\) 内的地块形状为 \(\triangle CDP\)，要求 \(A\)，\(B\) 均在线段 \(MN\) 上，\(C\)，\(D\) 均在圆弧上. 设 \(OC\) 与 \(MN\) 所成的角为 \(\theta\).

\begin{enumerate}
\item 用 \(\theta\) 分别表示矩形 \(ABCD\) 和 \(\triangle CDP\) 的面积，并确定 \(\sin\theta\) 的取值范围；
\item 若大棚 \(\mathrm{I}\) 内种植甲种蔬菜，大棚 \(\mathrm{II}\) 内种植乙种蔬菜，且甲，乙两种蔬菜的单位面积年产值之比为 \(4:3\)，求当 \(\theta\) 为何值时，能使甲，乙两种蔬菜的年总产值最大.
\end{enumerate}

\begin{center}
\bitmapfigure[width=0.58\linewidth]{img/2018/jiangsu/q17_fig2.png}
\end{center}
\end{problem}

\begin{solution}
\begin{enumerate}
\item 以 \(MN\) 所在直线为 \(x\) 轴，\(OP\) 所在垂线为 \(y\) 轴建立坐标系，则 \(P=(0,50)\)，\(O=(0,10)\)，\(MN\) 为 \(y=0\). 设 \(C=(40\cos\theta,10+40\sin\theta)\)（\(\cos\theta>0\)），\(D\) 与 \(C\) 关于 \(y\) 轴对称. 矩形宽为 \(80\cos\theta\)，高为 \(10+40\sin\theta\)，故面积为 \(800\left(4\sin\theta\cos\theta+\cos\theta\right)\text{ 平方米}\)；\(\triangle CDP\) 以 \(CD=80\cos\theta\) 为底，高为 \(50-(10+40\sin\theta)=40(1-\sin\theta)\)，故面积为 \(1600\left(\cos\theta-\sin\theta\cos\theta\right)\text{ 平方米}\). \(A\)，\(B\) 在线段 \(MN\) 上要求 \(40\cos\theta\leqslant OM\)，而 \(M\) 在圆上给出 \(OM=10\sqrt{15}\)，即 \(\cos\theta\leqslant\frac{\sqrt{15}}{4}\)，亦即 \(\sin\theta\geqslant\frac14\)；又 \(C\) 在 \(MN\) 上方且 \(\cos\theta>0\)，故 \(\sin\theta\in\left[\frac14,1\right)\).
\item 设甲、乙单位面积年产值分别为 \(4p\)，\(3p\)（\(p>0\)），则年总产值 \(Y=4p\cdot800\cos\theta(1+4\sin\theta)+3p\cdot1600\cos\theta(1-\sin\theta)=8000p\cos\theta(1+\sin\theta)\). 记 \(t=\sin\theta\in\left[\frac14,1\right)\)，\(Y^2\) 正比于 \(h(t)=(1-t)(1+t)^3\)，\(h'(t)=(1+t)^2(2-4t)\)，故 \(h\) 在 \(t=\frac12\) 处取得唯一极大值，即 \(\theta=\frac\pi6\) 时年总产值最大.
\end{enumerate}
\end{solution}

\begin{problem}
如图，在平面直角坐标系 \(xOy\) 中，椭圆 \(C\) 过点 \(\left(\sqrt3,\dfrac12\right)\)，焦点 \(F_1(-\sqrt3,0)\)，\(F_2(\sqrt3,0)\)，圆 \(O\) 的直径为 \(F_1F_2\).

\begin{enumerate}
\item 求椭圆 \(C\) 及圆 \(O\) 的方程；
\item 设直线 \(l\) 与圆 \(O\) 相切于第一象限内的点 \(P\).
\begin{enumerate}
\item 若直线 \(l\) 与椭圆 \(C\) 有且只有一个公共点，求点 \(P\) 的坐标；
\item 直线 \(l\) 与椭圆 \(C\) 交于 \(A\)，\(B\) 两点. 若 \(\triangle OAB\) 的面积为 \(\dfrac{2\sqrt6}{7}\)，求直线 \(l\) 的方程.
\end{enumerate}
\end{enumerate}

\begin{center}
\bitmapfigure[width=0.52\linewidth]{img/2018/jiangsu/q18_fig1.png}
\end{center}
\end{problem}

\begin{solution}
\begin{enumerate}
\item \(2a=|PF_1|+|PF_2|=\sqrt{(2\sqrt3)^2+\left(\frac12\right)^2}+\frac12=\frac72+\frac12=4\)，故 \(a=2\)；\(c=\sqrt3\)，\(b=\sqrt{a^2-c^2}=1\). 椭圆 \(C\) 的方程为 \(\frac{x^2}{4}+y^2=1\)，圆 \(O\) 的直径为 \(F_1F_2\)，方程为 \(x^2+y^2=3\).
\item
\begin{enumerate}
\item 设切点 \(P=(x_0,y_0)\)（\(x_0>0\)，\(y_0>0\)，\(x_0^2+y_0^2=3\)），则切线 \(l\) 为 \(x_0x+y_0y=3\). 代入椭圆得 \((y_0^2+4x_0^2)x^2-24x_0x+(36-4y_0^2)=0\). 有且仅有一个公共点当且仅当判别式为零，即 \(144x_0^2-(y_0^2+4x_0^2)(36-4y_0^2)=0\)；结合 \(y_0^2=3-x_0^2\) 化为 \((x_0^2-2)(x_0^2-3)=0\). \(x_0^2=3\) 给出 \(y_0=0\)，此时切线 \(x=\sqrt3\) 与椭圆交于两点，舍去；故 \(x_0^2=2\)，\(P=(\sqrt2,1)\).
\item 沿用上式记号，直线 \(l\) 与椭圆交于两点时弦长为 \(|AB|=\sqrt{1+\left(\frac{x_0}{y_0}\right)^2}|x_A-x_B|=\frac{4\sqrt{x_0^2-2}}{1+x_0^2}\). 原点到 \(l\) 的距离恒为 \(\frac{3}{\sqrt{x_0^2+y_0^2}}=\sqrt3\)，故 \(\triangle OAB\) 的面积为 \(\frac{2\sqrt3\sqrt{x_0^2-2}}{1+x_0^2}=\frac{2\sqrt6}{7}\)，即 \(\frac{\sqrt{x_0^2-2}}{1+x_0^2}=\frac{\sqrt2}{7}\). 记 \(u=x_0^2\) 得 \(2u^2-45u+100=0\)，解得 \(u=\frac52\)（\(u=20\) 因 \(u\leqslant3\) 舍去）. 于是 \(x_0=\frac{\sqrt{10}}{2}\)，\(y_0=\frac{\sqrt2}{2}\)，直线 \(l\) 的方程为 \(y=-\sqrt5x+3\sqrt2\).
\end{enumerate}
\end{enumerate}
\end{solution}

\begin{problem}
记 \(f'(x)\)，\(g'(x)\) 分别为函数 \(f(x)\)，\(g(x)\) 的导函数. 若存在 \(x_0\in\R\)，满足 \(f(x_0)=g(x_0)\) 且 \(f'(x_0)=g'(x_0)\)，则称 \(x_0\) 为函数 \(f(x)\) 与 \(g(x)\) 的一个“\(S\) 点”.

\begin{enumerate}
\item 证明：函数 \(f(x)=x\) 与 \(g(x)=x^2+2x-2\) 不存在“\(S\) 点”；
\item 若函数 \(f(x)=ax^2-1\) 与 \(g(x)=\ln x\) 存在“\(S\) 点”，求实数 \(a\) 的值；
\item 已知函数 \(f(x)=-x^2+a\)，\(g(x)=\dfrac{b\e^x}{x}\). 对任意 \(a>0\)，判断是否存在 \(b>0\)，使函数 \(f(x)\) 与 \(g(x)\) 在区间 \((0,+\infty)\) 内存在“\(S\) 点”，并说明理由.
\end{enumerate}
\end{problem}

\begin{solution}
\begin{enumerate}
\item 若 \(x_0\) 为“\(S\) 点”，则 \(1=f'(x_0)=g'(x_0)=2x_0+2\)，得 \(x_0=-\frac12\). 但 \(f\left(-\frac12\right)=-\frac12\neq-\frac{11}{4}=g\left(-\frac12\right)\)，故两函数不存在“\(S\) 点”.
\item 由 \(f'(x_0)=g'(x_0)\) 得 \(2ax_0=\frac1{x_0}\)（\(x_0>0\) 为 \(\ln x\) 的定义域要求），即 \(a=\frac1{2x_0^2}\)；代入 \(f(x_0)=g(x_0)\) 得 \(\frac12-1=\ln x_0\)，即 \(x_0=\frac1{\sqrt\e}\). 故存在“\(S\) 点”当且仅当 \(a=\frac\e2\).
\item \(f'(x_0)=-2x_0\)，\(g'(x)=\frac{b\e^x(x-1)}{x^2}\)（\(x>0\)）. 由 \(f'(x_0)=g'(x_0)\) 得 \(b=\frac{2x_0^3\e^{-x_0}}{1-x_0}\)，\(b>0\) 要求 \(x_0\in(0,1)\). 再由 \(f(x_0)=g(x_0)\) 得 \(a=-x_0^2+\frac{b\e^{x_0}}{x_0}=-x_0^2+\frac{2x_0^2}{1-x_0}=\frac{x_0^2(3-x_0)}{1-x_0}=:\varphi(x_0)\). \(\varphi\) 在 \((0,1)\) 上连续，\(\lim_{x_0\to0^+}\varphi(x_0)=0\)，\(\lim_{x_0\to1^-}\varphi(x_0)=+\infty\)，故对任意 \(a>0\) 存在 \(x_0\in(0,1)\) 使 \(\varphi(x_0)=a\)，相应 \(b=\frac{2x_0^3\e^{-x_0}}{1-x_0}>0\). 因此对任意 \(a>0\) 存在 \(b>0\) 使两函数在 \((0,+\infty)\) 内存在“\(S\) 点”.
\end{enumerate}
\end{solution}

\begin{problem}
设 \(\{a_n\}\) 是首项为 \(a_1\)，公差为 \(d\) 的等差数列，\(\{b_n\}\) 是首项为 \(b_1\)，公比为 \(q\) 的等比数列.

\begin{enumerate}
\item 设 \(a_1=0\)，\(b_1=1\)，\(q=2\)，若 \(|a_n-b_n|\le b_1\) 对 \(n=1\)，\(2\)，\(3\)，\(4\) 均成立，求 \(d\) 的取值范围；
\item 若 \(a_1=b_1>0\)，\(m\in\N^*\)，\(q\in(1,\sqrt[m]{2}]\)，证明：存在 \(d\in\R\)，使得 \(|a_n-b_n|\le b_1\) 对 \(n=2\)，\(3\)，\(\cdots\)，\(m+1\) 均成立，并求 \(d\) 的取值范围（用 \(b_1\)，\(m\)，\(q\) 表示）.
\end{enumerate}
\end{problem}

\begin{solution}
\begin{enumerate}
\item 此时 \(a_n=(n-1)d\)，\(b_n=2^{n-1}\). \(n=1\) 时 \(|a_1-b_1|=0\leqslant1\) 恒成立；\(n=2,3,4\) 分别给出 \(|d-2|\leqslant1\)，\(|2d-4|\leqslant1\)，\(|3d-8|\leqslant1\)，即 \(d\in[1,3]\)，\(d\in\left[\frac32,\frac52\right]\)，\(d\in\left[\frac73,3\right]\)，交集为 \(d\in\left[\frac73,\frac52\right]\).
\item 此时 \(a_n=b_1+(n-1)d\)，\(b_n=b_1q^{n-1}\). 条件即对 \(k=n-1=1,\ldots,m\) 有 \(|kd-b_1(q^k-1)|\leqslant b_1\)，亦即 \(d\in\left[\frac{b_1(q^k-2)}{k},\frac{b_1q^k}{k}\right]\). 由于 \(q^m\leqslant2\)，对 \(1\leqslant k\leqslant m\) 有 \(\frac{q^k}{k}\geqslant\frac{q^m}{m}\)（等价于 \(\frac{m}{m-j}\geqslant q^j\)（\(j=m-k\)）；而 \(q^j=(q^m)^{\frac{j}{m}}\leqslant2^{\frac{j}{m}}\)，函数 \((1-t)2^t\) 在 \([0,1)\) 上递减且初值为 \(1\)，故 \(\frac1{1-t}\geqslant2^t\) 即得）以及 \(\frac{q^k-2}{k}\leqslant\frac{q^m-2}{m}\)（等价于 \(mq^k-kq^m\leqslant2(m-k)\)；左端等于 \(jq^m-m(q^m-q^{m-j})\leqslant jq^m\leqslant2j\) 即得）. 故 \(m\) 个区间的交集恰为 \(\left[\frac{b_1(q^m-2)}{m},\frac{b_1q^m}{m}\right]\)，非空；存在性得证，\(d\) 的取值范围即为该区间.
\end{enumerate}
\end{solution}

\section{附加题：解答题}

\begin{problem}[A]
如图，圆 \(O\) 的半径为 \(2\)，\(AB\) 为圆 \(O\) 的直径，\(P\) 为 \(AB\) 延长线上一点，过 \(P\) 作圆 \(O\) 的切线，切点为 \(C\). 若 \(PC=2\sqrt3\)，求 \(BC\) 的长.

\begin{center}
\bitmapfigure[width=0.48\linewidth]{img/2018/jiangsu/q21_a_fig1.png}
\end{center}
\end{problem}

\begin{solution}
连 \(OC\)，则 \(OC\perp PC\). \(OP=\sqrt{OC^2+PC^2}=\sqrt{4+12}=4\)，而 \(OB=2\)，故 \(B\) 为斜边 \(OP\) 的中点，\(BC=\frac12OP=2\).
\end{solution}

\reuseproblemnumber
\begin{problem}[B]
已知矩阵 \(A=\begin{bmatrix}2&3\\1&2\end{bmatrix}\).

\begin{enumerate}
\item 求 \(A\) 的逆矩阵 \(A^{-1}\)；
\item 若点 \(P\) 在矩阵 \(A\) 对应的变换作用下得到点 \(P'(3,1)\)，求点 \(P\) 的坐标.
\end{enumerate}
\end{problem}

\begin{solution}
\begin{enumerate}
\item \(\det A=4-3=1\)，故 \(A^{-1}=\begin{bmatrix}2&-3\\-1&2\end{bmatrix}\).
\item 由 \(P=A^{-1}P'\) 得点 \(P\) 的坐标为 \((2\cdot3-3\cdot1,-1\cdot3+2\cdot1)=(3,-1)\).
\end{enumerate}
\end{solution}

\reuseproblemnumber
\begin{problem}[C]
在极坐标系中，直线 \(l\) 的方程为 \(\rho\sin\!\left(\dfrac{\pi}{6}-\theta\right)=2\)，曲线 \(C\) 的方程为 \(\rho=4\cos\theta\)，求直线 \(l\) 被曲线 \(C\) 截得的弦长.
\end{problem}

\begin{solution}
直线方程化为 \(\frac{x}{2}-\frac{\sqrt3y}{2}=2\) 即 \(x-\sqrt3y-4=0\)；曲线方程两边同乘 \(\rho\) 得 \(x^2+y^2=4x\)，即圆心 \((2,0)\)、半径 \(2\) 的圆. 圆心到直线的距离为 \(\frac{|2-4|}{\sqrt{1+3}}=1\)，故弦长为 \(2\sqrt{4-1}=2\sqrt3\).
\end{solution}

\reuseproblemnumber
\begin{problem}[D]
若 \(x\)，\(y\)，\(z\) 为实数，且 \(x+2y+2z=6\)，求 \(x^2+y^2+z^2\) 的最小值.
\end{problem}

\begin{solution}
由柯西不等式 \((1^2+2^2+2^2)(x^2+y^2+z^2)\geqslant(x+2y+2z)^2=36\)，得 \(x^2+y^2+z^2\geqslant4\)，等号在 \(\frac{x}{1}=\frac{y}{2}=\frac{z}{2}\)（即 \(x=\frac23\)，\(y=z=\frac43\)）时成立，故最小值为 \(4\).
\end{solution}



\begin{problem}
如图，在正三棱柱 \(ABC-A_1B_1C_1\) 中，\(AB=AA_1=2\)，点 \(P\)，\(Q\) 分别为 \(A_1B_1\)，\(BC\) 的中点.

\begin{enumerate}
\item 求异面直线 \(BP\) 与 \(AC_1\) 所成角的余弦值；
\item 求直线 \(CC_1\) 与平面 \(AQC_1\) 所成角的正弦值.
\end{enumerate}

\begin{center}
\bitmapfigure[width=0.56\linewidth]{img/2018/jiangsu/q22_fig1.png}
\end{center}
\end{problem}

\begin{solution}
\begin{enumerate}
\item 取 \(A=(0,0,0)\)，\(B=(2,0,0)\)，\(C=(1,\sqrt3,0)\)，\(A_1=(0,0,2)\)，则 \(B_1=(2,0,2)\)，\(P=(1,0,2)\)，\(Q=\left(\frac32,\frac{\sqrt3}{2},0\right)\)，\(\overrightarrow{BP}=(-1,0,2)\)，\(\overrightarrow{AC_1}=(1,\sqrt3,2)\). 异面直线所成角取锐角，余弦值为 \(\frac{|\overrightarrow{BP}\cdot\overrightarrow{AC_1}|}{|\overrightarrow{BP}||\overrightarrow{AC_1}|}=\frac{3}{\sqrt5\cdot\sqrt8}=\frac{3\sqrt{10}}{20}\).
\item \(\overrightarrow{AQ}=\left(\frac32,\frac{\sqrt3}{2},0\right)\)，平面 \(AQC_1\) 的法向量可取 \(\bs n=\overrightarrow{AQ}\times\overrightarrow{AC_1}=(\sqrt3,-3,\sqrt3)\)，\(|\bs n|=\sqrt{15}\). 直线 \(CC_1\) 方向为 \((0,0,1)\)，所成角的正弦值为 \(\frac{|\bs n\cdot(0,0,1)|}{|\bs n|}=\frac{\sqrt3}{\sqrt{15}}=\frac{\sqrt5}{5}\).
\end{enumerate}
\end{solution}

\begin{problem}
设 \(n\in\N^*\)，对 \(1\)，\(2\)，\(\cdots\)，\(n\) 的一个排列 \(i_1i_2\cdots i_n\)，如果当 \(s<t\) 时，有 \(i_s>i_t\)，则称 \((i_s,i_t)\) 是排列 \(i_1i_2\cdots i_n\) 的一个逆序，排列 \(i_1i_2\cdots i_n\) 的所有逆序的总个数称为其逆序数. 例如：对 \(1\)，\(2\)，\(3\) 的一个排列 \(231\)，只有两个逆序 \((2,1)\)，\((3,1)\)，则排列 \(231\) 的逆序数为 2. 记 \(f_n(k)\) 为 \(1\)，\(2\)，\(\cdots\)，\(n\) 的所有排列中逆序数为 \(k\) 的全部排列的个数.

\begin{enumerate}
\item 求 \(f_3(2)\)，\(f_4(2)\) 的值；
\item 求 \(f_n(2)\)（\(n\ge5\)）的表达式（用 \(n\) 表示）.
\end{enumerate}
\end{problem}

\begin{solution}
\begin{enumerate}
\item 逆序数为 \(2\) 的 \(3\) 元排列为 \(231\)，\(312\)，故 \(f_3(2)=2\)；\(4\) 元排列中逆序数为 \(2\) 的为 \(2314\)，\(3124\)，\(1342\)，\(1423\)，\(2143\)，故 \(f_4(2)=5\).
\item 记排列的逆序表为 \(l_s=\#\{t>s:i_t<i_s\}\)，则 \(\sum l_s=2\). 若某 \(l_s=2\)，则 \(i_s=s+2\)（更大则 \(l_s\geqslant3\)，更小则 \(l_s\leqslant1\)），且 \(l_{s+1}=0\) 迫使 \(s+1\)，\(s+2\) 位为 \(s\)，\(s+1\)，共 \(n-2\) 个. 若恰两个 \(l_s=l_t=1\)（\(s<t\)）：\(t\geqslant s+2\) 时 \(i_s=s+1\)，\(i_t=t+1\)，且 \(s\)，\(t\) 分居 \(s+1\)，\(t+1\) 位，即两个不交的相邻交换，共 \(\mathrm C_{n-1}^2-(n-2)=\frac{(n-2)(n-3)}{2}\) 个；\(t=s+1\) 时只能是三轮换 \(s+1\)，\(s+2\)，\(s\)，共 \(n-2\) 个. 总计 \(f_n(2)=\frac{(n-2)(n-3)}{2}+2(n-2)=\frac{(n-2)(n+1)}{2}=\frac{n^2-n-2}{2}\)（\(n\geqslant5\)）.
\end{enumerate}
\end{solution}
